%
% 16-reellanalytisch.tex -- Reelle differenzierbare Funktionen sind nicht alle analytisch
%
% (c) 2026 Prof Dr Andreas Müller
%
\subsection{Reelle differenzerierbare Funktionen sind nicht alle analytisch}
Die Aussage des Satzes~\ref{buch:analytisch:satz:holomorph-analytisch}
hat keine Entsprechung für reelle beliebig oft differenzierbare
Funktionen.
Vielmehr gibt es beliebig oft differenzierbare reellwertige
Funktionen auf $\mathbb{R}$, die keine konvergente Potenzreihe haben.
Wir konstruieren in diesem Abschnitt ein solches Beispiel.

%
% Eine glatte Funktion mit verschwindender Potenzreihe
%
\subsubsection{Eine glatte Funktion mit verschwindender Potenzreihe}
\input{chapters/040-analytisch/fig/fig-nichtanalytisch.tex}%
Wir betrachten die Funktion
\begin{equation}
f(x)
=
\begin{cases}
e^{-1/x}&\qquad \text{für $x>0$}\\
0&\qquad\text{sonst}
\end{cases}
\label{buch:analytisch:fig:nichtanalytisch}
\end{equation}
(siehe auch Abbildung~\ref{buch:analytisch:fig:nichtanalytisch}).
Sie ist beliebig oft differenzierbar.
Die Ableitung ist ausserhalb des Punktes $x=0$ stetig, es muss also
nur noch untersucht werden, dass die Ableitung auch im Punkt $x=0$
stetig ist.

%
% Die Ableitungen von e^{-1/x} im Punkt 0
%
\subsubsection{Die Ableitungen von $e^{-1/x}$ im Punkt $0$}
Die Ableitung von $f(x)$ ist
\begin{equation}
f'(x)
=
\frac{1}{x^2}
e^{-1/x}
\label{buch:analytisch:eqn:f'}
\end{equation}
für $x>0$.
Für höhere Ableitungen wird der Ausdruck komplizierter, wir wollen
mit vollständiger Induktion zeigen, dass alle Ableitungen von der Form
\begin{equation}
f^{(n)}(x)
=
P_n\biggl(\frac1x\biggr) e^{-1/x}
\label{buch:analytisch:eqn:fn=Pne}
\end{equation}
ist, wobei $P_n$ ein Polynom ist.
Die Gleichung \eqref{buch:analytisch:eqn:f'} dient dazu als
Induktionsverankerung.

Für den Induktionsschritt nehmen wir jetzt an, dass 
\eqref{buch:analytisch:eqn:fn=Pne}
für die $n$-te Ableitung gilt und zeigen, dass auch die $n+1$-te
Ableitung diese Form hat.
Dazu berechnen wir
\begin{align*}
f^{(n+1)}(x)
&=
\frac{d}{dx}
f^{(n)}(x)
=
\frac{d}{dx} P\biggl(\frac1x\biggr) e^{-1/x}
\\
&=
-P'_n\biggl(\frac1x\biggr) \frac{1}{x^2} e^{-1/x}
+
P_n\biggl(\frac1x\biggr) \frac{1}{x^2} e^{-1/x}
\\
&=
\frac{1}{x^2}
\biggl(
P_n\biggl(\frac1x\biggr)
-
P_n'\biggl(\frac1x\biggr)
\biggr)e^{-1/x}
\end{align*}
Setzt man 
\begin{equation}
P_{n+1}(X)
=
X^2(P_n(X) - P'_n(X)),
\label{buch:analytisch:eqn:e1rekursion}
\end{equation}
kann man tatsächlich
\[
f^{(n+1)}(x)
=
P_{n+1}\biggl(\frac1x\biggr)e^{-1/x}
\]
schreiben.
Die Rekursionsformel~\eqref{buch:analytisch:eqn:e1rekursion}
liefert ein Polynom $P_{n+1}$, womit der Induktionsschritt vollzogen
ist.

Die Polynome $P_n$ aus dem vorangegangenen Beweis lassen sich
mit der Rekursionsformel~\eqref{buch:analytisch:eqn:e1rekursion}
sehr leicht berechnen, die ersten paar Polynome sind
\begin{align*}
P_0(X) &= 1 \\
P_1(X) &= X^2 \\
P_2(X) &= X^4 - 2X^3 \\
P_3(X) &= X^2 (X^4 - 2X^3 - 4X^3 + 6X^2) = X^6 -6X^5+6X^4 \\
P_4(X) &= X^8 -12X^7+36X^6-24X^5 \\
P_5(X) &= X^{10} -20X^9 +120 X^8 - 240X^7 + 120 X^6.
\end{align*}
Da für $x\to 0$ der Exponentialfaktor $e^{-1/x}$ immer schneller
gegen $0$ geht, als ein polynomialer Faktor $P_n(1/x)$ gegen $\infty$
gehen kann, ist der Grenzwert 
\[
\lim_{x\to 0+}
f^{(n)}(x)
=
\lim_{\xi\to \infty}
P_n(\xi) e^{-\xi}
=
0.
\]
Damit ist gezeigt, dass an der Stelle $x=0$ alle Ableitungen von $f$
stetig sind, $f$ ist eine glatte Funktion.

Da die Ableitungen an der Stelle $x=0$ verschwinden, ist die
Taylor-Reihe von $f$ die Nullfunktion.
Da $f$ selbst nicht die Nullfunktion ist, kann die Taylor-Reihe
nicht gegen $f$ konvergieren.
%
% Die Funktion f(x) ist nicht Einschränkung einer holomorphen Funktion
%
\subsubsection{Die Funktion $f(x)$ ist nicht Einschränkung einer holomorphen Funktion}
Die Funktion $f(x)$ lässt sich nicht zu einer holomorphen Funktion
ausdehnen.
Der einzige dafür in Frage kommende Kandidat ist die Funktion
$z\mapsto e^{-1/z}$.
Nähert sich $z$ auf der imaginären Achse dem Punkt $0$, zum Beispiel
mit $z=i/t$ für $t\to \infty$, dann ist
\[
\lim_{\varepsilon\to 0} e^{-1/i\varepsilon}
=
\lim_{t\to\infty} e^{it}
=
\lim_{t\to\infty}(\cos t + i\sin t).
\]
Der Grenzwert auf der rechten Seite konvergiert allerdings nicht, sodass
die Funktion $z\mapsto e^{-1/z}$ im Punkt $z=0$ nicht einmal stetig
sein kann.

